
%\subsection{Vokabular}
%\beign{itemize}
%\item Spektrum
%\item Fourier-Theorie vs Potenzreihen
%\item Analyse und Synthese
%\item Schwingende Saite
%\item Partielle Differentialgleichung
%\item Diskrete Energiezust�nde
%\item Randbedingung
%\item Separationsl�sung
%\item Darstellung von Grundschwingung und Oberschwingungen: Beobachtung (Randbedingungen!)
%\item Wavelets
%\item Vektorraum, additive geschlossen
%\item Periode
%\item Bestapproximation mit trigonometrischen Funktionen
%\item Konvergenzverhalten der Fourier-Reihe
%\item Vollst�ndigkeit?
%	\end{itemize}

\ifintegriereaufgaben

\section*{Fourierreihen Aufgaben}\label{MA2-fourierreihen-aufgaben}
Erkl�rungen und Beispiele, die zur Bearbeitung der Aufgaben hilfreich sind, finden Sie im Skript\wennreferenzvorhanden{MA2-fourierreihen}{ (Abschnitt \abschnittsverweis{MA2-fourierreihen})} und, falls Sie nicht weiterkommen, dann schauen Sie \hyperref[MA2-fourierreihen-loesungen]{hier}.
\else
\section{Fourierreihen Aufgaben}\label{MA2-fourierreihen-aufgaben}

Regeln und Beispiele zu diesen Aufgaben finden Sie im Skript. Die L�sungen finden Sie \hyperref[MA2-fourierreihen-loesungen]{hier}.
\fi

	\begin{exercisebox}[Harmonische Schwingungen]
		Bestimmen Sie die Amplitude, die Periode und die Nullphase der folgenden harmonischen Schwingungen und skizzieren sie den Funktionsverlauf
		\begin{eqnarray*}
		(a)\; y(t) = 2\,\sin\left(3\, t-\dfrac{\pi}{6} \right)\quad \quad (b)\; y(t) = 10\,\sin\left(\pi\, t - 3\,\pi\right)\,.
		\end{eqnarray*}
\end{exercisebox}

	\begin{exercisebox}[komplexes Fourier-Polynom vs. reelles Fourier-Polynom]
		\begin{itemize}
			\item[(a)]
		Zeigen Sie, dass man die komplexwertigen Koffizienten $c_k$ eines komplexen Fourier Polynoms
		\[
			p_{N,f}(t) = \sum\limits_{k=-N}^N c_k \cdot e^{i\,k\,\omega_0\, t} \,,
		\]
		in die reell-wertigen Koeffizieinten $a_n$ und $b_n$ des entsprechenden reellen Fourier-Polynom umrechnen kann.
\item[(b)]
		Zeigen Sie, dass man die reellen Koffizienten eines Fourier Polynoms
		\[
			p_{N,f}(t) = \frac{a_0}{2} + \sum\limits_{n=1}^N\left( a_n \cdot \cos( n\, \omega_0\, t) + b_n \cdot \sin(n \, \omega_0 \, t)\right)\,,
		\]
		in die komplex-wertigen Koeffizienten $c_k$ des entsprechend komplex-wertigen Fourier-Polynoms umrechnen kann.
	\end{itemize}
\end{exercisebox}


\begin{exercisebox}[Fourier-Reihen]
	Skizzieren Sie zun�chst die Funktionen und berechnen Sie anschlie�end die Fourier-Reihen.
	\begin{itemize}
		\item[(a)] Periode $T=2$ und
			\[ f(t) = t, \quad t \in [0,2]\,.\]
		\item[(b)] Periode $T=1$ und
			\[ f(t) = |t|, \quad t \in \left[-\frac12,\frac12\right]\,.\]
		\item[(c)] Periode $T>0$ und
			\[ f(t) = t\cdot (T-t), \quad t \in [0,T]\,.\]
		\item[(d)] Periode $T=2\pi$ und
			\[ f(t) = \max(\sin(t),0), \quad t \in [0,2\pi]\,.\]

	\end{itemize}
\end{exercisebox}
